\documentclass[12pt]{article}
\usepackage[utf8]{inputenc}
\usepackage[T1]{fontenc}
\usepackage{hyperref}
\title{Embera Tad\'o Language Resources for K-12 Curriculum Development}
\author{Compiled by Claude Sonnet 5 (claude-sonnet-5) for Andr\'as Kornai}
\date{}
\begin{document}
\maketitle

\section{Introduction}
Embera Tad\'o (ISO 639-3 tdc; Glottocode embe1261) is a Chocoan language of Colombia. Search results gave no resources specific to the Tad\'o variety. Most material covers Embera in general, or other varieties such as Northern Embera or Alto Sin\'u (Kat\'io). Sources also disagree about where Tad\'o is spoken.

\section{Primary Educational Resources}
\begin{itemize}
\item A Universidad de C\'ordoba journal article notes that few pedagogical tools exist for Embera in schools, and discusses primers for the Alto Sin\'u variety, not Tad\'o. \url{https://revistas.unicordoba.edu.co/index.php/avedhum/article/download/3848/6138}
\end{itemize}

\section{Linguistic and Reference Materials}
\begin{itemize}
\item Ministry of Culture of Colombia, Embera document from the native-languages app series. \url{https://www.mincultura.gov.co:443/areas/poblaciones/APP-de-lenguas-nativas/Documents/Embera.pdf}
\item Gobernaci\'on de Antioquia, ethnolinguistic dictionary (PDF) that classifies Tad\'o among the southern Embera dialects. \url{https://antioquia.gov.co/images/PDF2/gerencia-indigena/Documentos_tecnicos_academicos/diccionario%20etnoling\%C3\%BC\%C3\%ADstico.pdf}
\item Northern Embera online dictionary (Webonary). This is a different variety. \url{https://www.webonary.work/northernembera/overview/introduction/}
\item Barre\~na Agirrebeitia and P\'erez-Caurel, ``The revitalization of Embera language in Colombia: from orality to literacy'' (Onomazein). \url{https://onomazein.letras.uc.cl/index.php/onom/article/view/30325}
\item UNIR repository record on an illustrated Embera dictionary as a teaching tool. I only saw it in the results and did not check its details. \url{https://reunir.unir.net/handle/123456789/14210}
\end{itemize}

\section{Religious and Community Resources}
None found.

\section{Digital and Audio Resources}
Only the Northern Embera online dictionary above, which is for a different variety.

\section{Curriculum Development Context}
Embera literacy work in Choc\'o began in the late twentieth century and has had education projects. Tad\'o-specific materials appear not to be available online. Starting from the wider Embera orthography and primers, with local Tad\'o consultation, is the plausible route.
\end{document}